\section{Closure and min-cut as comparison tools}\label{app:flow}

This appendix records the alternative closure oracle only to locate the
\FSC{} within the existing framework. Maximum closure admits
a minimum-cut representation \citep{Picard1976}; parametric maximum flow
and ratio-closure applications are treated by \citet{GalloEtAl1989}. The
graph-based predecessor of both is \citet{LerchsGrossmann1965}, whose
normalized tree is a forest with branch aggregates rather than a flow;
\citet{Hochbaum2001} recovers a flow from that tree in linear time and gives a
parametric variant with the cost of a single run, which makes it a comparison
tool on the ratio task as well.
These are comparison tools, not subroutines required by
\cref{alg:forest,alg:outer}.

For a fixed real parameter $\lambda$, assign vertex value
$a_i=p_i-\lambda w_i$. Add source s and sink t, capacities
\[
 c_{si}=\max\{a_i,0\},\qquad
 c_{it}=\max\{-a_i,0\},
\]
and a precedence-arc capacity $H>\sum_i\max\{a_i,0\}$.
Every minimum cut has a source-side vertex set C that is a closure, and its
capacity is
\[
 \sum_i\max\{a_i,0\}-[p(C)-\lambda w(C)].
\]
Thus minimizing the cut maximizes the penalized closure objective.
For a fixed $\lambda$, this construction works with either sign of the
parameter. A linear parametric implementation can use a different
normalization on a specified interval, as explained in the supplied manuscript.

At a breakpoint, the closure maximizing the penalized objective need not be
unique. In a residual network of a maximum flow, source-reachable vertices
give the \emph{minimal} source side. The complement of vertices that can
reach t gives the \emph{maximal} source side. A canonical comparison must use
the latter convention. For a ratio root, the empty closure can also tie at
value zero, so an implementation must preserve or recover a nonempty ratio
maximizer.

The planned experiments will compare methods on the same task: one ratio,
one capacity, or the full positive path. The worst-case bounds of parametric
flow cannot be used as a bound for the number of \FS{} pivots.
